\section{引言}

大模型应用的快速扩张，正把推理系统从单机实验环境推向面向生产的高并发服务平台。随着国产 AI 加速芯片、服务器整机与集群软件生态逐步成熟，围绕国产算力构建高性能、可维护、可扩展的推理引擎，已成为学界和产业界共同关注的方向。但“国产算力推理引擎”并不是把既有 serving 框架迁移到另一类设备上的简单重复。相较通用硬件环境，国产平台往往同时叠加设备能力离散、编译链约束差异、运行时治理复杂和交付路径多样等因素，使许多原本可以局部吸收的问题迅速上升为系统级问题。平台边界、运行时状态与工程闭环一旦同时变化，推理系统的设计重心也会随之迁移。与训练系统相比，推理引擎更关心动态批处理、显存复用、流式时延、长上下文，以及工具调用、多模态和结构化输出下的稳定性；对国产平台而言，还要额外处理硬件后端差异、算子完备性、编译运行时兼容性，以及上游框架快速演进带来的维护压力。

近年来，以 vLLM 为代表的推理框架通过 PagedAttention、连续批处理与高效内存管理显著改写了大模型服务系统的设计范式，SGLang 等系统则进一步强调请求编排、结构化控制流和复杂工作负载下的执行效率\citep{kwon2023efficient,vllm,sglang}。通用 serving 文献也表明，推理系统的研究重心正在持续外移：Orca 与 Sarathi-Serve 说明，迭代级调度、连续批处理和 chunked prefill 已成为高吞吐服务的基础组织方式\citep{orca2023,sarathi2023}；DistServe 与 Splitwise 表明，prefill 和 decode 的阶段解耦正在从局部优化上升为整体 goodput 设计问题\citep{zhong2024distserve,splitwise2024}；XGrammar 与 ElasticMM 则进一步把演进推向结构化执行与多模态并行，说明推理引擎正在从 token 生成器演变为状态密集型运行时\citep{dong2024xgrammar,liu2025elasticmm}。这意味着，当推理系统进入国产芯片、国产编译运行时和本土部署环境后，关键问题会迅速从单一硬件后端优化转向统一框架抽象与本土设备能力的持续对齐。此时，调度、缓存、多模态和结构化执行的可迁移性、可维护性与稳定性，往往比单点优化更重要。从产业侧看，国产算力推理引擎面向的也早已不只是学术基准测试，而是在线问答、代码助手、政企知识服务、具身智能中间层和 Agent 工作流等多种 AGI4S 场景；这些场景既要求较高的吞吐密度和资源利用率，也要求流式首 token 时延、上下文切换开销、错误恢复能力和服务可观测性达到生产标准。因此，国产推理引擎的竞争焦点正从“能否完成基础部署”转向“能否在真实业务下稳定、高效且可持续演进地运行”。

在这一背景下，开源主干复用增强路线、平台一体化方案和国产原生/独立路线虽然都在解决国产部署问题，但三者分化的根源并不只是优化手段不同，更在于对系统边界和演进节奏的假设不同：前者强调上游复用，并允许在共享主干外侧持续注入平台与场景增强，其中既包括直接围绕主线后端边界展开的适配，也包括以 vLLM-HUST 为代表的本土增强实践；中者强调平台统一性交付；后者则更侧重围绕自研状态机、缓存治理和部署闭环组织系统能力。相较训练侧更容易用 FLOPS、集群规模等指标描述阶段进展，推理侧的核心矛盾更多分布在请求形态、缓存复用、阶段干扰、故障恢复与升级节奏等难以被单一数字概括的维度上。尤其当模型族从纯文本扩展到多模态、语音和 Agent 工作流后，系统能否在较长时间尺度内维持接口兼容、性能稳定和运维可控，往往比一次局部 benchmark 的领先更能决定其工程价值。基于此，全文采用“一个底座、三条主路径”的分析骨架：底座是统一抽象、统一指标和可观测证据链，用来支撑多平台、多版本和多阶段优化条件下的可组合、可比较与可复现；三条主路径分别是执行与通信、状态治理、压缩与成本协同，对应执行骨架收敛、长上下文与高并发下的状态组织，以及单位 token 成本进入系统闭环。在此基础上，后文依次讨论生态格局与路线分化、分析骨架与核心架构、基于骨架的路线比较、由路线比较收束出的关键挑战，以及未来可能的收敛方向。全文重点不在于为某一路线背书，而在于识别哪些能力应沉淀为能力契约，哪些状态必须纳入运行时统一管理，哪些优化适合以插件或外围工具链独立演进，以及哪些 benchmark 足以支撑路线选择。

为使讨论建立在可交叉核验的材料之上，本文还对讨论范围和证据边界作如下限定。首先，本文的“国产算力推理引擎”既包括直接面向国产芯片的原生推理系统，也包括基于主流 serving 主干的国产后端适配、面向多类国产硬件的部署工具链，以及围绕服务状态机、缓存治理和交付闭环组织的独立运行时；但不把单纯的硬件整机、一体机交付页面或应用层封装直接等同于推理引擎本体。其次，本文优先使用四类证据：系统论文与技术报告用于界定机制与设计目标，官方文档与开源仓库用于确认能力边界和工程入口，标准或白皮书用于辅助描述生态收敛趋势，社区部署记录仅作为观察真实摩擦点的补充材料，而不单独支撑性能结论。再次，文中对路线比较主要围绕可演进性、状态治理能力、复杂负载承载能力和交付确定性展开，而不试图在缺少统一复现实验条件的前提下给出跨平台绝对性能排序。本文更关注复杂度被压在哪里、证据是否足以支撑相关判断，以及系统是否具备可持续演进边界，而不是把不同公开口径简单折算为单一优劣结论。

\subsection{研究对象与证据方法}

本文的综述对象，限定为近年在公开论文、技术报告、官方文档或开源仓库中能够形成交叉验证的国产算力推理相关系统与工程实践。就时间范围而言，文章重点覆盖大模型 serving 主干与国产后端适配快速演进的近年公开材料，并在必要处回溯少量更早的基础工作，用于交代连续批处理、PagedAttention、阶段解耦和结构化执行等机制来源。就对象纳入而言，只有当某一对象至少能够在“机制说明”“工程入口”“公开实现”三者中满足其中两项时，本文才将其纳入路线比较或章节主线；若某材料仅以新闻稿、发布会口径、单次演示或未公开复现链路的技术验证出现，则只作为生态线索或趋势性观察引用，不直接据此给出成熟度与性能判断。

基于这一原则，本文将证据分为四层使用。第一层是系统论文与技术报告，用于界定执行、状态、压缩和评测等机制问题的研究脉络；第二层是官方文档、源码仓库与安装部署入口，用于确认对象的实际能力边界、组件关系与可进入性；第三层是标准、白皮书与生态报告，用于辅助刻画国产生态的分工变化与标准化趋势；第四层才是社区部署记录和兼容性经验，它们仅用于补充真实工程摩擦点，而不单独支撑横向性能结论。按这一分层，正文中的比较表与路线判断优先回答三类问题：一是复杂度主要被压在何处，二是哪些系统边界已经具有较稳定的公开证据，三是哪些判断仍应保留为条件性的工程观察。这样做的目的，不是把本文伪装为严格元分析，而是把综述中的描述性事实、比较性判断与趋势性推断尽量分层组织，使不同性质的材料不被直接并列。
